\renewcommand{\xchaptertitle}{附录\thechapter}
\chapter{译注来源总表}
\begingroup\small\setstretch{1.02}
本表汇总全书 \textbf{144} 条学生向译注的来源与等级，按讲排列。每条译注在正文中以其概念首次出现处的脚注形式出现。「位置」一栏指该译注在源文件中的位置，记作 L\,$n$\,:\,$m$，即第 $n$ 讲源文件第 $m$ 行。

来源等级：\textbf{A} 为原书或作者文献、经典教材（带版次与章号）；\textbf{B} 为 nLab、Stacks Project、MathWorld、DLMF、EMS 等正式参考源；\textbf{I} 为本库既往译注与术语表；等级标作「译者按」者表示无外部依据、仅据上下文推断。所有网页来源的访问日期为 2026-11-14。

\begin{longtable}{@{}p{1.0cm}p{3.3cm}p{1.5cm}p{1.3cm}p{7.9cm}@{}}
\toprule
\textbf{讲} & \textbf{概念} & \textbf{位置} & \textbf{等级} & \textbf{来源} \\
\midrule
\endfirsthead
\toprule
\textbf{讲} & \textbf{概念} & \textbf{位置} & \textbf{等级} & \textbf{来源} \\
\midrule
\endhead
\bottomrule
\endfoot
1 & 完美域 & \texttt{L01:22} & B & \url{https://ncatlab.org/nlab/show/perfect+field} \\
1 & Krull拓扑 & \texttt{L01:37} & B & \url{https://ncatlab.org/nlab/show/profinite+group} \\
1 & 可分闭包 & \texttt{L01:45} & B & \url{https://ncatlab.org/nlab/show/separable+closure} \\
1 & 阿贝尔化 & \texttt{L01:74} & B & \url{https://ncatlab.org/nlab/show/abelianization} \\
1 & 分裂域 & \texttt{L01:79} & B & \url{https://ncatlab.org/nlab/show/splitting+field} \\
1 & Dirichlet L-级数 & \texttt{L01:139} & B & \url{https://ncatlab.org/nlab/show/L-function} \\
1 & 逆Galois问题 & \texttt{L01:141} & B & \url{https://ncatlab.org/nlab/show/inverse+Galois+problem} \\
1 & 射有限完备化 & \texttt{L01:176} & B & \url{https://ncatlab.org/nlab/show/profinite+group} \\
\midrule
2 & 理想类群 & \texttt{L02:4} & B & \url{https://ncatlab.org/nlab/show/ideal+class+group} \\
2 & 代数整数环 & \texttt{L02:6} & A & Neukirch, Algebraic Number Theory, Grundlehren der mathematischen Wissenschaften 322, 第 I 章 §8（本讲参考文献） \\
2 & 分歧指数 & \texttt{L02:13} & A & Neukirch, Algebraic Number Theory, Grundlehren der mathematischen Wissenschaften 322, 第 I 章 §8（本讲参考文献） \\
2 & 剩余域 & \texttt{L02:15} & A & Neukirch, Algebraic Number Theory, Grundlehren der mathematischen Wissenschaften 322, 第 I 章 §8（本讲参考文献） \\
2 & 完全分裂 & \texttt{L02:19} & A & Neukirch, Algebraic Number Theory, Grundlehren der mathematischen Wissenschaften 322, 第 I 章 §9（本讲参考文献） \\
2 & 分解群 & \texttt{L02:67} & A & Neukirch, Algebraic Number Theory, Grundlehren der mathematischen Wissenschaften 322, 第 I 章 §9（本讲参考文献） \\
2 & 惯性群 & \texttt{L02:81} & A & Neukirch, Algebraic Number Theory, Grundlehren der mathematischen Wissenschaften 322, 第 I 章 §9（本讲参考文献） \\
2 & 互反映射 & \texttt{L02:111} & A & Cassels--Fröhlich, Algebraic Number Theory, 第 VII 章（本讲参考文献） \\
2 & 素位 & \texttt{L02:147} & B & \url{https://ncatlab.org/nlab/show/place} \\
2 & 射线类群 & \texttt{L02:159} & A & Neukirch, Algebraic Number Theory, Grundlehren der mathematischen Wissenschaften 322, 第 VI 章（本讲参考文献） \\
\midrule
3 & 局部整体原理 & \texttt{L03:9} & A & Neukirch, Algebraic Number Theory, Grundlehren der mathematischen Wissenschaften 322, 第 VI 章（本讲参考文献） \\
3 & 整体域 & \texttt{L03:18} & A & Neukirch, Algebraic Number Theory, Grundlehren der mathematischen Wissenschaften 322, 第 VI 章（本讲参考文献） \\
3 & 局部域 & \texttt{L03:23} & A & Neukirch, Algebraic Number Theory, Grundlehren der mathematischen Wissenschaften 322, 第 V 章（本讲参考文献） \\
3 & 范数子群 & \texttt{L03:71} & A & Neukirch, Algebraic Number Theory, Grundlehren der mathematischen Wissenschaften 322, 第 V 章（本讲参考文献） \\
3 & 阿代尔环 & \texttt{L03:92} & A & Cassels–Fröhlich, Algebraic Number Theory, 第 II 章（本讲参考文献） \\
3 & 限制直积 & \texttt{L03:94} & A & Cassels–Fröhlich, Algebraic Number Theory, 第 II 章（本讲参考文献） \\
3 & 伊代尔群 & \texttt{L03:96} & A & Cassels–Fröhlich, Algebraic Number Theory, 第 II 章（本讲参考文献） \\
\midrule
4 & 自守表示 & \texttt{L04:4} & A & Bernstein–Gelbart, An introduction to the Langlands Program, Birkhäuser 2004, 第 7 章（本讲参考文献） \\
4 & 同余子群 & \texttt{L04:148} & A & Bernstein–Gelbart, 同上, 第 7 章（本讲参考文献） \\
4 & 自守形式 & \texttt{L04:162} & B & \url{https://ncatlab.org/nlab/show/automorphic+form} \\
4 & 适度增长 & \texttt{L04:176} & A & Bernstein–Gelbart, 同上, 第 7 章（本讲参考文献） \\
4 & 抛物子群 & \texttt{L04:187} & B & \url{https://ncatlab.org/nlab/show/parabolic+subgroup} \\
4 & 中心特征标 & \texttt{L04:193} & A & Bernstein–Gelbart, 同上, 第 7 章（本讲参考文献） \\
4 & 尖点自守形式 & \texttt{L04:193} & A & Bernstein–Gelbart, 同上, 第 7 章（本讲参考文献） \\
\midrule
5 & (g,K)-模（正文作 (g,O)-模） & \texttt{L05:11} & A & Bernstein–Gelbart, An introduction to the Langlands Program, 第 7 章（本讲参考文献） \\
5 & 可容许表示 & \texttt{L05:28} & A & Bernstein–Gelbart, 同上, 第 7 章（本讲参考文献） \\
5 & 限制张量积 & \texttt{L05:55} & A & Bernstein–Gelbart, 同上, 第 7 章（本讲参考文献） \\
5 & 超尖点表示 & \texttt{L05:119} & A & Bernstein–Gelbart, 同上, 第 7 章（本讲参考文献） \\
5 & 主系列表示 & \texttt{L05:119} & A & Bernstein–Gelbart, 同上, 第 7 章（本讲参考文献） \\
5 & Satake参数 & \texttt{L05:140} & A & Bernstein–Gelbart, 同上, 第 7 章（本讲参考文献） \\
\midrule
6 & 模形式 & \texttt{L06:8} & B & \url{https://ncatlab.org/nlab/show/modular+form} \\
6 & 同余子群 & \texttt{L06:14} & B & \url{https://ncatlab.org/nlab/show/congruence+subgroup} \\
6 & Hecke算子 & \texttt{L06:16} & A & D. Bump, Automorphic Forms and Representations, Cambridge Studies in Advanced Mathematics 55, Cambridge University Press, 2004（本讲参考文献） \\
6 & q展开 & \texttt{L06:47} & B & \url{https://ncatlab.org/nlab/show/modular+form} \\
6 & Eisenstein级数 & \texttt{L06:60} & B & \url{https://ncatlab.org/nlab/show/Eisenstein+series} \\
6 & Eisenstein 展开的验算 & \texttt{L06:64} & C & 代入 $k=2,3$ 验证系数 240、-504；若用 $B_k$ 则 $k=2$ 得 48。上下文推断，无外部依据。 \\
6 & 新形式 & \texttt{L06:89} & C & 上下文推断，无外部依据 \\
\midrule
7 & 椭圆曲线 & \texttt{L07:8} & B & \url{https://ncatlab.org/nlab/show/elliptic+curve} \\
7 & j不变量 & \texttt{L07:29} & B & \url{https://ncatlab.org/nlab/show/j-invariant} \\
7 & Tate模 & \texttt{L07:76} & B & \url{https://ncatlab.org/nlab/show/Tate+module} \\
7 & 同源 & \texttt{L07:78} & A & J. H. Silverman, The arithmetic of elliptic curves, Graduate Texts in Mathematics 106, Springer, 2009（本讲参考文献） \\
7 & $\ell$进表示 & \texttt{L07:95} & A & J. H. Silverman, The arithmetic of elliptic curves, Graduate Texts in Mathematics 106, Springer, 2009（本讲参考文献）；并见 https://ncatlab.org/nlab/show/Galois+representation \\
7 & 超奇异/普通 & \texttt{L07:116} & B & \url{https://ncatlab.org/nlab/show/elliptic+curve} \\
7 & 良约化/坏约化 & \texttt{L07:120} & A & J. H. Silverman, The arithmetic of elliptic curves, Graduate Texts in Mathematics 106, Springer, 2009（本讲参考文献） \\
\midrule
8 & 紧化 & \texttt{L08:13} & B & \url{https://ncatlab.org/nlab/show/compactification}（Last revised 2025-08-19） \\
8 & 模曲线 & \texttt{L08:30} & B & \url{https://ncatlab.org/nlab/show/modular+curve}（Last revised 2015-07-21） \\
8 & 模空间 & \texttt{L08:87} & B & \url{https://ncatlab.org/nlab/show/moduli+space}（Last revised 2025-10-21） \\
8 & Hecke代数 & \texttt{L08:101} & A & D. Diamond, J. Shurman, A First course in Modular forms, GTM 228 (2005), Section 7.9 \\
8 & 双陪集 & \texttt{L08:126} & B & \url{https://ncatlab.org/nlab/show/double+coset}（Last revised 2019-08-01） \\
8 & Jacobian簇 & \texttt{L08:170} & B & \url{https://ncatlab.org/nlab/show/Jacobian+variety}（Last revised 2026-09-26） \\
8 & 新形式 & \texttt{L08:207} & A & D. Diamond, J. Shurman, A First course in Modular forms, GTM 228 (2005), Chapter 9 \\
\midrule
9 & 函数域 & \texttt{L09:72} & B & \url{https://ncatlab.org/nlab/show/function+field}（Last revised 2022-06-05） \\
9 & 素位 & \texttt{L09:94} & B & \url{https://ncatlab.org/nlab/show/place}（Last revised 2015-12-17） \\
9 & 阿代尔环 & \texttt{L09:114} & B & \url{https://ncatlab.org/nlab/show/ring+of+adeles}（Last revised 2025-07-01） \\
9 & 局部系统 & \texttt{L09:153} & B & \url{https://ncatlab.org/nlab/show/local+system}（Last revised 2025-04-10） \\
9 & 原书 p79 定理 2.4 未写出的条件（标准表述，译者按） & \texttt{L09:166} & A & 本讲参考文献 [D] Drinfeld, Amer. J. Math. 105 (1983) 85-114；[L] Lafforgue, Invent. Math. 147 (2002) 1-241 \\
\midrule
10 & 局部系统 & \texttt{L10:4} & B & \url{https://ncatlab.org/nlab/show/local+system} \\
10 & 平展基本群 & \texttt{L10:4} & B & \url{https://ncatlab.org/nlab/show/fundamental+group}（另见 Milne, Lectures on etale cohomology） \\
10 & 有限平展覆盖 & \texttt{L10:26} & B & \url{https://ncatlab.org/nlab/show/etale+morphism+of+schemes} \\
10 & 定理1.4原书未写Hausdorff & \texttt{L10:35} & 译者按 & 无外部来源（znoteplain，上下文推断） \\
10 & 引理2.2原书缺正规化有限与正规曲线两假设 & \texttt{L10:107} & 译者按 & 无外部来源（znoteplain，上下文推断） \\
10 & 光滑层 & \texttt{L10:124} & B & \url{https://ncatlab.org/nlab/show/constructible+sheaf}（条目较短，另见 Milne, Lectures on etale cohomology） \\
10 & 单值表示 & \texttt{L10:146} & 译者按 & 无外部来源（znoteplain，上下文推断；原引 nLab "monodromy" 不涉 Out 记号，按来源抽查降级） \\
\midrule
11 & Tate层 & \texttt{L11:27} & B & \url{https://ncatlab.org/nlab/show/Tate+twist} \\
11 & 两条代数事实原书缺Noether假设 & \texttt{L11:88} & 译者按 & 无外部来源（znoteplain，上下文推断） \\
11 & 向量丛 & \texttt{L11:110} & B & \url{https://ncatlab.org/nlab/show/vector+bundle} \\
11 & 模叠 & \texttt{L11:116} & B & \url{https://ncatlab.org/nlab/show/moduli+stack} \\
11 & 下降数据 & \texttt{L11:140} & B & \url{https://ncatlab.org/nlab/show/descent} \\
\midrule
12 & Hecke对应 & \texttt{L12:50} & B & \url{https://ncatlab.org/nlab/show/Hecke+correspondence} \\
12 & 几何Langlands & \texttt{L12:107} & B & \url{https://ncatlab.org/nlab/show/Langlands+program} \\
12 & Hecke特征层 & \texttt{L12:113} & B & \url{https://ncatlab.org/nlab/show/Hecke+eigensheaf} \\
\midrule
13 & Verdier对偶 & \texttt{L13:4} & B & \url{https://ncatlab.org/nlab/show/Verdier+duality}（另见 Peters–Steenbrink, Mixed Hodge Structures, Ch. 13 §13.1） \\
13 & 反常层 & \texttt{L13:115} & B & \url{https://ncatlab.org/nlab/show/perverse+sheaf} \\
13 & t结构 & \texttt{L13:141} & B & \url{https://ncatlab.org/nlab/show/t-structure} \\
13 & 反常上同调 & \texttt{L13:200} & B & \url{https://ncatlab.org/nlab/show/perverse+sheaf} \\
13 & 左右t-正合是一对约定 & \texttt{L13:211} & 译者按 & 无外部来源（nLab "t-structure" 无 left/right exact 定义段；按细则用 znoteplain，未引 BBD 定理号） \\
13 & 交复形 & \texttt{L13:215} & B & \url{https://ncatlab.org/nlab/show/intersection+cohomology} \\
13 & Artin指子对象降链条件（非阿廷叠） & \texttt{L13:247} & 译者按 & 无外部来源（znoteplain，上下文推断） \\
\midrule
14 & Riemann–Hilbert对应 & \texttt{L14:2} & B & \url{https://ncatlab.org/nlab/show/Riemann-Hilbert+correspondence} \\
14 & D模 & \texttt{L14:4} & A & J. Bernstein, Algebraic Theory of D-modules（本讲参考文献 [B]）；并见 https://ncatlab.org/nlab/show/D-module \\
14 & 凝聚D模 & \texttt{L14:31} & A & Peters–Steenbrink, Mixed Hodge Structures, Ch. 13 §13.3–§13.6（本讲参考文献；原并引 nLab "coherent D-module" 仅 582 字符且无定义，按来源抽查删除） \\
14 & 平坦联络 & \texttt{L14:47} & B & \url{https://ncatlab.org/nlab/show/flat+connection} \\
14 & 原书 p117 的 (f·ξ)(w)=ξ(f·w) 易错提示（可见译注） & \texttt{L14:56} & 译者按 & 上下文推断，无外部依据（原书 p117 表述含混，见报告校订清单 \#5 后说明） \\
14 & 内射对象 & \texttt{L14:71} & B & \url{https://ncatlab.org/nlab/show/injective+object} \\
14 & 转移双模与 f\_+ 的方向易错点（可见译注） & \texttt{L14:101} & 译者按 & 上下文推断，无外部依据（原书 p119 未给出显式写法；按来源抽查删去译者自补的显式公式与双模记号，只保留方向提示） \\
14 & 特征簇 & \texttt{L14:163} & B & \url{https://ncatlab.org/nlab/show/holonomic+D-module}；并见 Peters–Steenbrink, Mixed Hodge Structures, Ch. 13 §13.3–§13.6 \\
14 & 完整D模 & \texttt{L14:180} & B & \url{https://ncatlab.org/nlab/show/holonomic+D-module}；并见 J. Bernstein, Algebraic Theory of D-modules \\
\midrule
15 & 刚性化 & \texttt{L15:10} & 译者按 & 上下文推断，无外部依据（nLab 无 rigidification 条目） \\
15 & 对称幂 & \texttt{L15:36} & B & \url{https://mathworld.wolfram.com/SymmetricPower.html} \\
15 & Abel–Jacobi映射 & \texttt{L15:38} & B & \url{https://ncatlab.org/nlab/show/Abel-Jacobi+map} \\
15 & 几何点 & \texttt{L15:104} & 译者按 & 上下文推断，无外部依据（原书混用几何点 $x$ 与 $k$-有理点 $p$） \\
15 & 射有限完备化 & \texttt{L15:132} & B & \url{https://ncatlab.org/nlab/show/profinite+group} \\
\midrule
16 & Artin映射 & \texttt{L16:18} & A & J. Neukirch, \textit{Algebraic Number Theory}, Grundlehren der mathematischen Wissenschaften 322，第 VI 章（见第 2 讲参考文献 [N]；来源库 work/nlab/alt/reciprocity\_map.md） \\
16 & Artin互反律 & \texttt{L16:20} & B & \url{https://ncatlab.org/nlab/show/Artin+reciprocity+law} \\
16 & 特征标层 & \texttt{L16:44} & 译者按 & 上下文推断，无外部依据（nLab character sheaf 条目 NOTFOUND，见 work/nlab/list\_l16.tsv） \\
16 & Lang同源 & \texttt{L16:56} & 译者按 & 上下文推断，无外部依据（nLab Lang isogeny 条目 NOTFOUND，见 work/nlab/list\_l16.tsv） \\
\midrule
17 & 带平坦联络的丛 & \texttt{L17:13} & B & \url{https://ncatlab.org/nlab/show/flat+connection} \\
17 & 阿贝尔簇 & \texttt{L17:78} & B & \url{https://ncatlab.org/nlab/show/abelian+variety} \\
17 & 对偶环面（必注） & \texttt{L17:89} & B & \url{https://ncatlab.org/nlab/show/Cartier+duality} \\
17 & 主极化 & \texttt{L17:89} & B & \url{https://mathworld.wolfram.com/AbelianVarietyPolarization.html} \\
17 & Fourier–Mukai变换（必注） & \texttt{L17:99} & B & \url{https://ncatlab.org/nlab/show/Fourier-Mukai+transform} \\
17 & 摩天层 & \texttt{L17:117} & B & \url{https://ncatlab.org/nlab/show/skyscraper+sheaf} \\
17 & θ 与经典 theta 除子相差一个逆（命题 3.2 后译者按） & \texttt{L17:149} & 译者按 & 原书 p146 未定义 S(L\_0)、S\_A(L\_0) 中取逆的约定，属约定不明；按 Lead 裁决照原文直译，丰沛方向以原书为准。 \\
\midrule
18 & 根数据（必注） & \texttt{L18:5} & A & Springer, Reductive groups, Proc. Symp. Pure Math. 33, AMS 1979, pp. 3-27（本讲参考文献 [S]）；nLab 无 root datum 条目（root+datum 404） \\
18 & 根系（必注） & \texttt{L18:17} & A & Springer, Reductive groups, Proc. Symp. Pure Math. 33, AMS 1979, pp. 3-27（本讲参考文献 [S]）；nLab 仅有 98 字符存根，不引 \\
18 & 约化群（必注） & \texttt{L18:54} & B & \url{https://ncatlab.org/nlab/show/reductive+group}（另参 Springer 1979, pp. 3-27） \\
18 & 极大环面 & \texttt{L18:70} & B & \url{https://ncatlab.org/nlab/show/maximal+torus} \\
18 & 相对根系 & \texttt{L18:118} & A & Springer, Reductive groups, Proc. Symp. Pure Math. 33, AMS 1979, pp. 3-27（本讲参考文献 [S]）；nLab 无相对根系条目 \\
18 & 对偶群（必注） & \texttt{L18:124} & A & Springer, Reductive groups, Proc. Symp. Pure Math. 33, AMS 1979, pp. 3-27（本讲参考文献 [S]）；nLab Langlands dual group 仅薄页，不引 \\
\midrule
19 & Bruhat分解（必注） & \texttt{L19:13} & B & \url{https://ncatlab.org/nlab/show/Bruhat+decomposition} \\
19 & 旗簇 & \texttt{L19:28} & B & \url{https://ncatlab.org/nlab/show/flag+variety} \\
19 & 球面Hecke代数（必注） & \texttt{L19:57} & A & I. Satake, Theory of spherical functions on reductive algebraic groups over p-adic fields, IHES Publ. Math. 18 (1963) 5-69（本讲参考文献 [S]）；nLab Hecke algebra 仅 98 字符存根，不引 \\
19 & 支配余权（必注） & \texttt{L19:69} & A & I. Satake, IHES Publ. Math. 18 (1963) 5-69（本讲参考文献 [S]）；nLab weight 条目仅 83 字符存根，不引 \\
19 & Satake同构（必注） & \texttt{L19:116} & A & I. Satake, IHES Publ. Math. 18 (1963) 5-69（本讲参考文献 [S]）；nLab 无代数 Satake 同构条目（仅有 geometric Satake equivalence，不引作代数侧来源） \\
19 & 带状球函数 & \texttt{L19:129} & B & \url{https://ncatlab.org/nlab/show/spherical+function} \\
\midrule
20 & 局部Langlands对应（自选） & \texttt{L20:106} & B & S. Gelbart, An elementary introduction to the Langlands program, Bull. Amer. Math. Soc. 10 (1984) 177-219（本讲参考文献 [G]） \\
20 & Langlands函子性（必注） & \texttt{L20:108} & B & \url{https://ncatlab.org/nlab/show/Langlands+functoriality} \\
20 & 整体Hecke代数的分解（译者按） & \texttt{L20:115} & 译者按 & 原书 p169 写 $\otimes$'\_v H(G(F\_v),GL\_n(O\_v))；数域情形可按「有限部分 K\_f + 无穷处相应代数」拆分，说明见正文 可见译注 \\
20 & L函数（自选） & \texttt{L20:125} & B & S. Gelbart, An elementary introduction to the Langlands program, Bull. Amer. Math. Soc. 10 (1984) 177-219（本讲参考文献 [G]） \\
20 & 未分歧表示（必注） & \texttt{L20:142} & A & J. Bernstein, S. Gelbart (eds.), An introduction to the Langlands Program, Birkhäuser 2004, 第 11 章（J. W. Cogdell）（本讲参考文献 [BG]） \\
20 & L群（必注） & \texttt{L20:144} & A & J. Bernstein, S. Gelbart (eds.), An introduction to the Langlands Program, Birkhäuser 2004, 第 11 章（J. W. Cogdell）（本讲参考文献 [BG]） \\
\midrule
21 & 仿射Grassmannian & \texttt{L21:15} & B & \url{https://ncatlab.org/nlab/show/affine+Grassmannian} \\
21 & 几何Satake等价 & \texttt{L21:17} & B & \url{https://ncatlab.org/nlab/show/geometric+Satake+equivalence} \\
21 & ind概形 & \texttt{L21:25} & B & \url{https://ncatlab.org/nlab/show/ind-scheme} \\
21 & 环群 & \texttt{L21:66} & B & \url{https://ncatlab.org/nlab/show/loop+group} \\
21 & 中心扩张 & \texttt{L21:68} & B & \url{https://ncatlab.org/nlab/show/group+extension} \\
21 & Kac–Moody代数 & \texttt{L21:70} & B & \url{https://ncatlab.org/nlab/show/affine+Lie+algebra} \\
21 & Cartan分解 & \texttt{L21:92} & 译者按 & - \\
\midrule
22 & 卷积积 & \texttt{L22:35} & B & \url{https://ncatlab.org/nlab/show/geometric+Satake+equivalence} \\
22 & 等变层 & \texttt{L22:41} & B & \url{https://ncatlab.org/nlab/show/equivariant+sheaf} \\
22 & 半小映射 & \texttt{L22:68} & A & de Cataldo–Migliorini, The decomposition theorem, perverse sheaves and the topology of algebraic maps, §4.2（本讲参考文献） \\
22 & 幺半范畴 & \texttt{L22:189} & B & \url{https://ncatlab.org/nlab/show/monoidal+category} \\
\midrule
23 & 张量函子 & \texttt{L23:18} & B & \url{https://ncatlab.org/nlab/show/monoidal+functor} \\
23 & 奇偶层 & \texttt{L23:78} & A & Juteau–Mautner–Williamson, Parity sheaves, §2.2（本讲参考文献） \\
23 & 半单性 & \texttt{L23:97} & B & \url{https://ncatlab.org/nlab/show/semisimple+category} \\
23 & Tannaka对偶 & \texttt{L23:120} & B & \url{https://ncatlab.org/nlab/show/Tannaka+duality} \\
\end{longtable}
\endgroup

\chapter{原书笔误与存疑登记}
\begingroup\small\setstretch{1.02}
下表登记在翻译与校订过程中发现的\textbf{英文原稿笔误}，以及已按正确写法处理的说明。按本工程纪律，原稿笔误一律不悄悄改动正文而不留痕：凡译文已按正确写法处理的，在对应 TeX 源码中以行注释（Tutor Note）就地登记。附表为其汇总，供读者与后续校订者核对，共 147 条。

\begin{longtable}{@{}p{1.0cm}p{2.2cm}p{10.6cm}@{}}
\toprule
\textbf{讲} & \textbf{位置} & \textbf{说明} \\
\midrule
\endfirsthead
\toprule
\textbf{讲} & \textbf{位置} & \textbf{说明} \\
\midrule
\endhead
\bottomrule
\endfoot
1 & \texttt{L01:87} & 原书 p6：theta 的余域前后不一致 \\
1 & \texttt{L01:114} & 原书 p7：推论 2.5 次数公式重复 zeta\_n \\
1 & \texttt{L01:127} & 原书 p8：定理 2.6 证明下界写错 \\
1 & \texttt{L01:131} & 原书 p8：order m\_1 旧译误作指数 \\
1 & \texttt{L01:137} & 原书 p10：Dirichlet 定理未声明 b != 0 \\
1 & \texttt{L01:149} & 原书 p8：习题提示漏 m=3 \\
2 & \texttt{L02:55} & 原书 p13：Example 1.2 第四个变量印作 c \\
2 & \texttt{L02:101} & 原书 p14：Remark 2.7(2) 把 x\^{}q 印作 x\^{}f \\
2 & \texttt{L02:113} & 原书 p15：Artin 映射定义域印作 I\^{}S\_K \\
2 & \texttt{L02:187} & 原书 p17：误引 Corollary 3.5 \\
3 & \texttt{L03:29} & 原书 p20 分解群条件 \\
3 & \texttt{L03:39} & 原书 p21 [L:K] \\
3 & \texttt{L03:44} & 原书 p21 v\_K 定义域 \\
3 & \texttt{L03:81} & 原书 p22 F\_p 记号 \\
3 & \texttt{L03:90} & 原书 p23 O\_p=R\_{$\ge$0} \\
3 & \texttt{L03:129} & 原书 p24 Gal(L/K) \\
3 & \texttt{L03:142} & 原书 p25 Theorem 3.3 \\
4 & \texttt{L04:5} & 原书 p27 措辞收窄 \\
4 & \texttt{L04:8} & 原书 p27 K 的宿主 \\
4 & \texttt{L04:13} & 原书 p28 指数缺负号 \\
4 & \texttt{L04:127} & 原书 p32 K 的宿主 \\
4 & \texttt{L04:142} & 原书 p32 定理 2.1 条件 \\
4 & \texttt{L04:147} & 原书 p33 G(R) 应为 G(Q\_S) \\
4 & \texttt{L04:179} & 原书 p34 Z(A) 为对角矩阵 \\
4 & \texttt{L04:190} & 原书 p35 抛物子群须真 \\
5 & \texttt{L05:7} & 原书 p36 G(A) 应为 G(A\^{}f) \\
5 & \texttt{L05:15} & 原书 p36 a\_0 应为 o\_0 \\
5 & \texttt{L05:100} & 原书 p39 ω\_{-1/2} 参数 \\
5 & \texttt{L05:122} & 原书 p40 μ=(μ1,μ1) \\
5 & \texttt{L05:125} & 原书 p40 (A\^{}$\times$)\^{}2 \\
5 & \texttt{L05:145} & 原书 p41 K\_0(r) 等三处 \\
5 & \texttt{L05:156} & 原书 p41 Remark 3.7 \\
5 & \texttt{L05:169} & 原书 p42 Casimir 与 Z(g) \\
6 & \texttt{L06:25} & 原书 p44：$B(\Q)$ 误作 $G(\R)$ 中上三角子群 \\
6 & \texttt{L06:65} & 原书 p46：Eisenstein 展开的下标与符号 \\
6 & \texttt{L06:88} & 原书 p47：指数误作 $2k$ \\
6 & \texttt{L06:102} & 原书 p47：系数估计的指数符号 \\
6 & \texttt{L06:141} & 原书 p49：Casimir 作用式的记号与符号 \\
6 & \texttt{L06:147} & 原书 p49：尖点积分漏 $j(g,i)^{-1}$ \\
6 & \texttt{L06:178} & 原书 p50：Dirichlet 特征标定义不全 \\
6 & \texttt{L06:194} & 原书 p51：$c(\pi_p)=1$ 与非分歧矛盾 \\
6 & \texttt{L06:200} & 原书 p51："oddforms" \\
6 & \texttt{L06:221} & 原书 p52：文献书名拼写 "Automophic" \\
7 & \texttt{L07:23} & 原书 p54：群法则把 P=Q 的情形写成 P+P=$\infty$ \\
7 & \texttt{L07:42} & 原书 p54：注 1.3 反例的换元 \\
7 & \texttt{L07:140} & 原书 p58：把 p 次幂态射称作 geometric Frobenius \\
7 & \texttt{L07:142} & 原书 p58：$f(t)$ 的括号位置 \\
7 & \texttt{L07:144} & 原书 p58：特征值所在域写成 $\mathbb{Z}_p$ \\
7 & \texttt{L07:154} & 原书 p58：$Z(E_p,p^{-s})$ 与 $Z(E_p,t)$ 不自洽 \\
7 & \texttt{L07:182} & 原书 p60：Euler 乘积漏特征标 \\
8 & \texttt{L08:77} & 原书 p64：$\omega\cong\Omega^1_{Y_\Gamma}$ 应为 $\omega^{\otimes2}\cong\Omega^1_{Y_\Gamma}$ \\
8 & \texttt{L08:109} & 原书 p66：式 (2.1) 第一个乘积下标 $p\mid N$ 应为 $p\nmid N$ \\
8 & \texttt{L08:134} & 原书 p67：$T_pf$ 式中 $\sum_{p\mid N}$ 应为 $\sum_{p\mid n}$ \\
8 & \texttt{L08:214} & 原书 p70：Hecke 代数记作 $T_k$，应为 $H_k$ \\
8 & \texttt{L08:238} & 原书 p71：作者名 "J. Churman" 应为 J. Shurman（DS7） \\
8 & \texttt{L08:243} & 原书 p71：作者名 "M. Papoport" 应为 M. Rapoport \\
8 & \texttt{L08:248} & 原书 p71：作者名 "J. Churman" 应为 J. Shurman（DS9） \\
9 & \texttt{L09:36} & 原书 p73：由 $1=|a_0(t)^n|\leq|u(t)|^n+R$ 直接断言 $|u(t)|=1$ \\
9 & \texttt{L09:125} & 原书 p77：中心写成 "the diagonal matrices" \\
9 & \texttt{L09:176} & 原书 p79：作者名 "V. G. Drinfel" 应为 Drinfeld \\
10 & \texttt{L10:109} & 引理2.2 正文回退原书（原书未提正规化有限与正规曲线） \\
10 & \texttt{L10:25} & 原书笔误：例1.2 的 A$\times$ 应为 C$\times$ \\
10 & \texttt{L10:147} & 原书笔误：Out 定义分母应为 Inn \\
10 & \texttt{L10:166} & 原书笔误：X\_F 应为 X\_k̄（译文已作 X\_k̄） \\
10 & \texttt{L10:60} & 原书排版笔误：ob(E) 应为 ob(C) \\
10 & \texttt{L10:144} & 原书交叉引用失效：(??) 应为 (3.1) \\
11 & \texttt{L11:13} & 原书笔误：连续映射 is zero 应为有限像 \\
11 & \texttt{L11:44} & 原书笔误：Aut(Z$\ell$(r)) 与 Diag 写法（应为 Z$\ell$$\times$ 与 χ$\ell$\^{}r$\in$GL1） \\
11 & \texttt{L11:127} & 原书术语错：lax functors 应为伪函子 \\
11 & \texttt{L11:143} & 原书下标错乱：下降数据余循环条件（译文用标准式） \\
11 & \texttt{L11:156} & 原书笔误：Hom(T,T) 应为 Hom(T,T') \\
12 & \texttt{L12:82} & 原书口径：Q̄$\ell$ 层写作 Z$\ell$ 层$\otimes$Q$\ell$，译文改用有限扩张 E \\
13 & \texttt{L13:207} & 原书笔误：additive 应作三角函子 \\
14 & \texttt{L14:58} & 原书 p117 右 D 模作用 d$\langle$ξ,w$\rangle$ 缺负号 \\
14 & \texttt{L14:79} & 原书 p118 de Rham 复形末项 Ω\_X\^{}{d\_X-1} \\
14 & \texttt{L14:102} & 原书 p119 的 f\^{}! 平移 [d\_Y-d\_X] \\
14 & \texttt{L14:107} & 原书 p119 的 μ\_Z(D\_Z) 下标 \\
14 & \texttt{L14:196} & 原书 p122 的 D\_X$\circ$f\_+ 下标 \\
14 & \texttt{L14:214} & 原书 p123 的 φ\_u 下标 q \\
14 & \texttt{L14:220} & 原书 p123 称 Sol(F) 为 D\_X-模 \\
16 & \texttt{L16:29} & 原书笔误：p135 绝对 Frobenius 的记号 \\
16 & \texttt{L16:55} & 原书笔误：p135 $\mathcal L_{\overline x}\otimes\mathcal O_{\overline y}$ \\
16 & \texttt{L16:67} & 原书笔误：p136 $\overline\rho_m$、$\overline\rho_\boxtimes$ 缺指数 $-1$ \\
17 & \texttt{L17:12} & 原书 p140：归纳定义末项作 $\nabla$\^{}s \\
17 & \texttt{L17:26} & 原书 p141：复形写成 C-代数层 \\
17 & \texttt{L17:35} & 原书 p141：命题编号 1.8 应为 1.3 \\
17 & \texttt{L17:140} & 原书 p145：命题 3.1 作 Φ\_{S(L)} \\
17 & \texttt{L17:152} & 原书 p146：命题 3.2 的 θ、θ' 取值 \\
17 & \texttt{L17:159} & 原书 p146：注记 3.3 作 θ\^{}{-1} 丰沛 \\
18 & \texttt{L18:15} & 原书 p148：V=Q$\otimes$Q 漏下标 Z \\
18 & \texttt{L18:29} & 原书 p149：V\_R 误作 V$\otimes$\_Q C \\
18 & \texttt{L18:48} & 原书 p150：嵌入 GL(k) 漏下标 n \\
18 & \texttt{L18:67} & 原书 p151：伴随矩阵对角元次序 \\
18 & \texttt{L18:75} & 原书 p151：T\_α maps identically \\
18 & \texttt{L18:78} & 原书 p152：T=(Im α$\vee$)$\times$T\_α 与配对缺 =2 \\
18 & \texttt{L18:83} & 原书 p152：例 2.3 作 X$\cong$Z\^{}2 \\
18 & \texttt{L18:99} & 原书 p152：(1)(b) 限制下标 V\_α \\
18 & \texttt{L18:111} & 原书 p153：约化 k 群用 k\^{}s，译文用 bar k \\
18 & \texttt{L18:121} & 原书 p154：绝对根系写成 Φ(G(k),S(k)) \\
18 & \texttt{L18:131} & 原书 p154：支配权 λ$\in$X$\vee$ 且配对次序 \\
18 & \texttt{L18:145} & 原书 p155：f(W)$\ne$0 应为 f(v)$\ne$0 \\
18 & \texttt{L18:148} & 原书 p155：non-constant 与正次数齐次 \\
18 & \texttt{L18:157} & 原书 p155：定义 3.1 漏「光滑」 \\
19 & \texttt{L19:8} & 原书 p157：U\_α 的 Lie 代数作 g\_α \\
19 & \texttt{L19:9} & 原书 p157–158：P\_θ 称为 P 的子群 \\
19 & \texttt{L19:16} & 原书 p158：n\_w\^{}{-1}U\_α n\_w=U\_{w(α)} \\
19 & \texttt{L19:23} & 原书 p158：U'\_w、U''\_w 的条件 \\
19 & \texttt{L19:25} & 原书 p158：(1.3) 右端漏 P(k) \\
19 & \texttt{L19:35} & 原书 p158：Bruhat 判据作「前缀」 \\
19 & \texttt{L19:64} & 原书 p160：引理 2.1 证明中 ξ(t) 记号 \\
19 & \texttt{L19:65} & 原书 p160：t$\in$O\_F\^{}$\times$ 与 t\_1=…=t\_r=1 \\
19 & \texttt{L19:71} & 原书 p160：Cartan 分解对 Z\^{}r\_{$\ge$0} 求和 \\
19 & \texttt{L19:114} & 原书 p162：γ\_i 为 simple coroots \\
19 & \texttt{L19:124} & 原书 p162：注记 3.2 的 x\_i(α)=e\^{}{$\langle$α,γ\_i$\rangle$} \\
19 & \texttt{L19:144} & 原书 p163：f\_χ(khu)=χ(h) 未归一化 \\
19 & \texttt{L19:153} & 原书 p163：X$\subset$X$'$={…} \\
20 & \texttt{L20:2} & 原书 p165：local field 未排除阿基米德情形 \\
20 & \texttt{L20:15} & 原书 p165：(1.1) 分母作求和 \\
20 & \texttt{L20:32} & 原书 p166：(1.2) 用 e\^{}{$-$ν}、e\^{}{$-$w(ν)} \\
20 & \texttt{L20:34} & 原书 p166：U(F) 说成 G\^{}L 正根所定 Borel \\
20 & \texttt{L20:57} & 原书 p167：(1.3) 用 e\^{}{$-$ν} \\
20 & \texttt{L20:67} & 原书 p167：取 ν=$-$λ 而本处取 ν=w\_0λ \\
20 & \texttt{L20:89} & 原书 p168：引 Definition 3.4（应为 3.3） \\
20 & \texttt{L20:97} & 原书 p168：注记 2.1 的 e\^{}{b\_j} 与「T 的一参数子环面」 \\
20 & \texttt{L20:146} & 原书 p171：L 同态未要求 G\^{}L\_0 上代数 \\
20 & \texttt{L20:160} & 原书 p172：猜想 2.5 漏写 f \\
20 & \texttt{L20:165} & 原书 p172：注记 2.6 作「put G'={1}」 \\
21 & \texttt{L21:9} & 原书 p173：Hecke 特征层误作 over G \\
21 & \texttt{L21:34} & 原书 p175：证明中 β\^{}{-1}$\circ$α\^{}{-1} 复合次序错 \\
21 & \texttt{L21:64} & 原书 p176：式 (2.1) 未标约化 \\
21 & \texttt{L21:84} & 原书 p177：正根定义哑标重复用 k \\
21 & \texttt{L21:89} & 原书 p177：伴随作用公式漏 t\^{}m 因子 \\
21 & \texttt{L21:105} & 原书 p178：P\_λ 误称 standard parabolic \\
21 & \texttt{L21:121} & 原书 p178：Remark 2.3(2) 纤维切空间漏环面方向 \\
21 & \texttt{L21:124} & 原书 p178：T(O) 分解误作等式 \\
21 & \texttt{L21:141} & 原书 p179：嵌入应取 H\^{}0 的对偶 \\
21 & \texttt{L21:150} & 原书 p179：定理 2.5(1) 首句漏 overline \\
21 & \texttt{L21:151} & 原书 p179：定理 2.5(2) 维数符号颠倒 \\
21 & \texttt{L21:168} & 原书 p179：P\_S 定义未提有限维支撑 \\
22 & \texttt{L22:53} & 原书 p184：反常性判据第二条件用 H\^{}j\_c 且符号与中围条件不符 \\
22 & \texttt{L22:78} & 原书 p185：半小映射不等式右端用像维数 dim f(T1) \\
22 & \texttt{L22:89} & 原书 p186：小映射定义 U/W 混用 \\
22 & \texttt{L22:122} & 原书 p187：定义1.6 条件1 与条件2、3 不等价 \\
23 & \texttt{L23:11} & 原书 p191：Remark 1.1 第二条件符号写作 t\^{}γ t\^{}{-s} \\
23 & \texttt{L23:47} & 原书 p193：τ° 定义后未写 [-2] 平移 \\
23 & \texttt{L23:83} & 原书 p195：P\_λ 作 standard parabolic \\
23 & \texttt{L23:142} & 原书 p197：用单余根集 Δ$\vee$ 生成支配余权 \\
\end{longtable}
\endgroup
